\chapter{Introduction to Weighted Directional Total Nuclear Variation}
\label{chapter:wdtnv}
%\SA{Check comment}
In this chapter, we introduce a smoothed \acrfull{wdtnv} prior for synergistic emission tomography. The chapter focuses on the regulariser itself: its relationship to total nuclear variation, its differentiability after singular-value smoothing, the use of anatomical information through directional operators, and the modality-wise weighting needed when coupling heterogeneous emission images.

The optimisation machinery used to apply this prior in large-scale \gls{pet}/\gls{spect} reconstruction, including prior-aware preconditioning, subset selection, and multi-bed \gls{pet} handling, is developed separately in Chapter~\ref{chapter:optimisation}.

\section{Total Nuclear Variation Revisited}
\label{sec:tv_to_tnv}

This section builds on the definitions in Section~\ref{subsec:synergistic_recon}. For convenience, we restate the key definitions here:
\begin{itemize}
    \item Registered images: $\bm{z} \coloneqq \{ z_m \}_{m=1}^{M}$, where $z_m \coloneqq W_m x_m$ and
    $W_m : \mathbb{R}_+^{J_m} \to \mathbb{R}_+^{N}$.
    \item Finite-difference operator: for $\ell=1,\dots,d$,
    \[
        \big((D x)_n\big)^{(\ell)} \coloneqq \frac{x_{n+\delta_\ell} - x_{n}}{\Delta_{\ell}} .
    \]
    \item Voxel-wise Jacobian:
    \[
        \mathcal{G}_n(\bm{z})
        \coloneqq
        \begin{bmatrix}
        D z_{1,n} & D z_{2,n} & \cdots & D z_{M,n}
        \end{bmatrix}
        \in \mathbb{R}^{d\times M}.
    \]
    \item Nuclear norm: for a matrix $Y$ with singular values $\{\sigma_k(Y)\}_{k=1}^{r}$ and $r=\operatorname{rank}(Y)$,
    \[
        \|Y\|_* \coloneqq \sum_{k=1}^{r}\sigma_k(Y).
    \]
    \item Total nuclear variation:
    \[
        \mathcal{R}_{\mathrm{TNV}}(\bm{z}) \coloneqq \sum_{n=1}^N  \|\mathcal{G}_n(\bm{z})\|_* .
    \]
\end{itemize}

\subsection{Mathematical properties and differentiability}
\label{subsec:tnv_differentiability}

Briefly revisiting total variation, optimisation of $\mathcal{R}_{\mathrm{TV}}$ (cf.\ \eqref{eq:tv}) is inconvenient in gradient-based methods because these functionals are not globally smooth: the TV gradient is undefined when $(D x_m)_j=0$. A standard remedy is to introduce a small smoothing parameter $\epsilon>0$.

If we define
\begin{equation}
\tilde{\mathcal{R}}_{\mathrm{TV}}(x_m)
\;:=\;
\sum_{j=1}^{J_m} \|(D x_m)_j\|_{2,\epsilon},
\qquad
\|v\|_{2,\epsilon}:=\sqrt{\|v\|_2^2+\epsilon^2},
\end{equation}
then the image-space gradient is
\begin{equation}
\nabla_{x_m}\tilde{\mathcal{R}}_{\mathrm{TV}}(x_m)
\;=\;
-\,\operatorname{div}\!\left(\frac{D x_m}{\|D x_m\|_{2,\epsilon}}\right),
\end{equation}
where the division is understood pointwise over voxels.

For nuclear-norm vectorial total variation, the unsmoothed singular-value map is not differentiable at zero singular values, and hence the corresponding nuclear norm is not $C^1$ at rank-deficient points. We therefore apply the Charbonnier function to the singular values:
\begin{equation}\label{eq:charbonnier_g}
g(s)=\sqrt{s^2+\epsilon^2}-\epsilon,
\qquad
g'(s) = \frac{s}{\sqrt{s^2+\epsilon^2}},
\qquad
g''(s)=\frac{\epsilon^2}{(s^2+\epsilon^2)^{3/2}}.
\end{equation}
The constant $-\epsilon$ matches the implemented objective value and does not affect gradients or optimisation.

The nuclear norm is orthogonally invariant: for any orthogonal $U\in\mathbb{R}^{d\times d}$ and $V\in\mathbb{R}^{M\times M}$,
\begin{equation}
\|Y\|_* \;=\; \|U Y V^\top\|_*.
\end{equation}
More generally, any function of the form
\begin{equation}
\phi(Y) \;:=\; f(\sigma(Y)),
\qquad
f(s)=\sum_{\ell=1}^{r} g(s_\ell),
\qquad
r=\min\{d,M\},
\end{equation}
is orthogonally invariant and depends on $Y$ only through its singular values $\sigma(Y)=(\sigma_1(Y),\dots,\sigma_r(Y))$.
Since $g\in C^\infty([0,\infty))$ for $\epsilon>0$, the symmetric function $f$ is smooth and hence $\phi$ is differentiable, with gradient given by the standard formula for orthogonally invariant matrix functions \cite{lewisConvexAnalysisUnitarily1995,lewisDerivativesSpectralFunctions1996}.

\paragraph{Jacobian-space gradient.}
Let $\phi(\mathcal{G})=\sum_{\ell=1}^{\min(d,M)} g(\sigma_\ell(\mathcal{G}))$, where $\sigma_\ell(\mathcal{G})$ are the singular values of the voxel-wise Jacobian matrix $\mathcal{G}\in\mathbb{R}^{d\times M}$. A thin \gls{svd}
\begin{equation}
\mathcal{G} \;=\; U\,\Sigma\,V^\top,
\qquad
U\in\mathbb{R}^{d\times r},\ \Sigma=\mathrm{diag}(\sigma)\in\mathbb{R}^{r\times r},\ V\in\mathbb{R}^{M\times r},
\end{equation}
with $r=\min\{d,M\}$ gives
\begin{equation}\label{eq:phi_grad_svd}
\nabla_{\mathcal{G}} \phi(\mathcal{G})
\;=\;
U\,\mathrm{diag}\!\bigl(g'(\sigma)\bigr)\,V^\top,
\qquad
g'(\sigma_\ell)=\frac{\sigma_\ell}{\sqrt{\sigma_\ell^2+\epsilon^2}}.
\end{equation}

In practice, it is unnecessary to form $U$ if we exploit the small Gram matrix
\begin{equation}
G := \mathcal{G}^\top \mathcal{G}\in\mathbb{R}^{M\times M},
\qquad G\succeq 0.
\end{equation}
Let $G=V\Lambda V^\top$ with $\Lambda=\mathrm{diag}(\lambda_1,\dots,\lambda_M)$. Then $\lambda_\ell=\sigma_\ell^2$ for the nonzero spectrum and $V$ coincides with the right singular vectors of $\mathcal{G}$. On the range of $\mathcal{G}$ we have $U=\mathcal{G}V\Sigma^{-1}$, so substituting into~\eqref{eq:phi_grad_svd} yields
\begin{equation}\label{eq:phi_grad_gram}
\nabla_{\mathcal{G}} \phi(\mathcal{G})
\;=\;
\mathcal{G}\,V\,\mathrm{diag}\!\left(\frac{g'(\sigma_\ell)}{\sigma_\ell}\right)\,V^\top,
\end{equation}
with the continuous extension at $\sigma_\ell=0$. For Charbonnier smoothing,
\begin{equation}
\frac{g'(\sigma_\ell)}{\sigma_\ell}
=
\frac{1}{\sqrt{\sigma_\ell^2+\epsilon^2}}
=
\frac{1}{\sqrt{\lambda_\ell+\epsilon^2}},
\end{equation}
hence
\begin{equation}\label{eq:phi_grad_inv_sqrt}
\nabla_{\mathcal{G}} \phi(\mathcal{G})
\;=\;
\mathcal{G}\,V\,\mathrm{diag}\!\left(\frac{1}{\sqrt{\lambda_\ell+\epsilon^2}}\right)\,V^\top
\;=\;
\mathcal{G}\,(G+\epsilon^2 \mathbb{I})^{-1/2}.
\end{equation}
Equation~\eqref{eq:phi_grad_inv_sqrt} shows that the Jacobian-space gradient is obtained by applying the inverse square-root of the $M\times M$ \gls{spd} matrix $G+\epsilon^2 \mathbb{I}$ to the columns of $\mathcal{G}$.

For $M=2$ modalities, even an eigendecomposition of $G$ can be avoided. Writing
\begin{equation}
G=
\begin{pmatrix}
a & c\\
c & b
\end{pmatrix},
\qquad
a=\|\mathcal{G}_{:,1}\|_2^2,\ \ b=\|\mathcal{G}_{:,2}\|_2^2,\ \ c=\langle \mathcal{G}_{:,1},\mathcal{G}_{:,2}\rangle,
\end{equation}
define $S:=G+\epsilon^2 \mathbb{I}=\begin{psmallmatrix}s_{11} & s_{12}\\ s_{12} & s_{22}\end{psmallmatrix}$ with $s_{11}=a+\epsilon^2$, $s_{22}=b+\epsilon^2$, $s_{12}=c$, and the invariants
\begin{equation}
t:=\mathrm{tr}(S)=s_{11}+s_{22},
\qquad
s:=\sqrt{\det(S)}=\sqrt{s_{11}s_{22}-s_{12}^2}.
\end{equation}
Since $\epsilon>0$ implies $S\succ 0$, the following identity holds for $2\times2$ \glsentryshort{spd} matrices:
\begin{equation}\label{eq:inv_sqrt_2x2}
S^{-1/2}
=
\sqrt{t+2s}\,(S+s\mathbb{I})^{-1}
=
\frac{\sqrt{t+2s}}{2\det(S)+s\,t}
\begin{pmatrix}
s_{22}+s & -s_{12}\\
-s_{12} & s_{11}+s
\end{pmatrix}.
\end{equation}
Therefore, for two modalities the Jacobian-space gradient can be evaluated via
\begin{equation}\label{eq:2x2_gram}
\nabla_{\mathcal{G}} \phi(\mathcal{G})
=
\mathcal{G}\,S^{-1/2},
\qquad S=G+\epsilon^2 \mathbb{I},
\end{equation}
using only dot products and scalar square-roots per voxel.

\paragraph{Image-space gradient.}
Define voxel-wise
\begin{equation}
Q_n \;:=\; \nabla_{\mathcal{G}} \phi\bigl(\mathcal{G}_n(\bm{z})\bigr)
\in\mathbb{R}^{d\times M},
\end{equation}
and let $Q=\{Q_n\}_{n=1}^N$. Applying the chain rule through the linear map $\bm{z}\mapsto D\bm{z}$ yields
\begin{equation}
\delta \tilde{\mathcal{R}}_{\mathrm{TNV}}(\bm{z})[\delta \bm{z}]
=
\sum_{n=1}^N \left\langle Q_n,\ (D\,\delta \bm{z})_n \right\rangle_F
=
\left\langle Q,\ D\,\delta \bm{z} \right\rangle.
\end{equation}
By definition of the adjoint $D^\ast$,
\begin{equation}
\left\langle Q,\ D\,\delta \bm{z} \right\rangle
=
\left\langle D^\ast Q,\ \delta \bm{z} \right\rangle.
\end{equation}
Defining the discrete divergence as $\operatorname{div}:=-D^\ast$, we obtain
\begin{equation}
\nabla_{\bm{z}} \tilde{\mathcal{R}}_{\mathrm{TNV}}(\bm{z})
=
-\,\operatorname{div} Q,
\end{equation}
where $\operatorname{div}$ acts column-wise on $Q$.

\section{Weighted Directional Total Nuclear Variation ($\wdtnv$)}
\label{sec:wdtnv}

This section extends TNV \cite{rigieJointReconstructionMultichannel2015, holtTotalNuclearVariation2014} in two orthogonal ways:
\begin{itemize}
    \item high-resolution anatomical information is incorporated directionally rather than as an additional coupled intensity channel;
    \item voxel-wise, modality-wise diagonal weighting compensates for heterogeneous scaling and spatially varying reliability across modalities.
\end{itemize}

The starting point is the standard TNV construction: penalising the nuclear norm of a multi-modality Jacobian promotes joint sparsity and alignment of edges, since a low nuclear norm favours a locally low-rank Jacobian.

\subsection{Directional operators from guidance}
\label{subsec:directional_from_guidance}

To incorporate guidance information, such as a high-resolution anatomical image, we borrow the operator from directional total variation (Equation~\ref{eq:directional_op}). The key design choice is that the guidance image is used only to steer the penalty, not to add an extra column to the joint Jacobian. This avoids directly coupling an anatomical channel, with distinct units, resolution, and artefacts, to the emission channels through the singular spectrum.

If a guidance image were appended as an additional modality, then a single global scaling would have data-dependent and spatially non-uniform consequences. Where emission gradients are weak, even modest guidance gradients can dominate the local Jacobian and hence the singular values; where emission gradients are strong, the same scaling may become negligible. Since the nuclear norm couples modalities through the eigenstructure of the local Gram matrix, such reweighting can also transfer structured guidance artefacts into the coupled reconstruction. By contrast, a directional operator modifies which gradient components are penalised, without forcing the presence of edges where the emission data do not support them.

Applying $\Pi_{m,n}$ to each modality gradient defines the directional multi-modality Jacobian
\begin{equation}
\mathcal{G}^{\mathrm{dir}}_n(\bm{z})= \bigl[\Pi_{1,n}(D z_1)_n \;|\; \cdots \;|\; \Pi_{M,n}(D z_M)_n\bigr]\in\mathbb{R}^{d\times M}.
\end{equation}
Here $D$ denotes the chosen discrete directional-difference operator and each $\Pi_{m,n}\in\mathbb{R}^{d\times d}$ is constructed from the guidance field at voxel $n$. In the implementation used in this work, a guidance image $u$ defines
\begin{equation}
\xi_n
=
\frac{(D u)_n}{\sqrt{\|(D u)_n\|_2^2+\eta^2}},
\qquad
\Pi_{m,n}=\mathbb{I}_d-\gamma_m\,\xi_n\xi_n^\top,
\qquad \gamma_m\in[0,1].
\label{eq:wdtnv_directional_operator}
\end{equation}
When $\eta$ is not specified explicitly it is set from a robust scale estimate of $\|(Du)_n\|_2$, namely $0.01$ times the difference between its $99$th and $1$st percentiles. In homogeneous regions of the guidance image, $\xi_n$ is small and $\Pi_{m,n}$ is close to the identity, so the regulariser reduces to an approximately isotropic TNV-type coupling.

Using \gls{ct} only to construct the directional operators decouples anatomical steering from emission-to-emission weighting. Modality weights determine how strongly emission gradients are coupled to one another through the nuclear norm, while $\Pi_{m,n}$ determines how strongly the penalty prefers guidance-consistent edge directions. Retaining a two-modality Jacobian also preserves the efficient $2\times2$ Gram-matrix structure derived in Section~\ref{sec:tv_to_tnv}.

\subsection{Modality-wise weighting and the full \texorpdfstring{$\wdtnv$}{w-dTNV} prior}
\label{subsec:modality_weighting}
\label{subsec:voxelwise_weighting}

\gls{pet} and \gls{spect} images of the same administered \gls{y90} activity differ substantially in dynamic range: the effective count level, sensitivity, and noise characteristics vary by orders of magnitude between the two modalities, and in clinical imaging the two reconstructions may also operate at different spatial resolutions. If the directional emission gradients $\Pi_{m,n}(D z_m)_n$ are used directly without normalisation, the nuclear-norm coupling is dominated by the modality with the larger gradient magnitudes.

To compensate, we introduce a voxel-wise, modality-wise diagonal weight matrix
\begin{equation}
    B_n := \mathrm{diag}(\rho_{1,n},\,\rho_{2,n}) \in \mathbb{R}^{2\times 2},
    \qquad \rho_{m,n} > 0.
    \label{eq:Bn_def}
\end{equation}
The weighted directional Jacobian is
\begin{equation}
    \tilde{\mathcal{G}}_n(\bm{z}) := \mathcal{G}_n^{\mathrm{dir}}(\bm{z})\,B_n
    = \bigl[\Pi_{1,n}(D z_1)_n \;\big|\; \Pi_{2,n}(D z_2)_n\bigr]\,
    \mathrm{diag}(\rho_{1,n},\rho_{2,n})
    \in \mathbb{R}^{d\times 2}.
    \label{eq:Yn_weighted}
\end{equation}
The columns of $\tilde{\mathcal{G}}_n$ are $\tilde{g}_{m,n} := \rho_{m,n}\,\Pi_{m,n}(D z_m)_n \in \mathbb{R}^d$.

The weighted directional Total Nuclear Variation ($\wdtnv$) prior is then
\begin{equation}
    \mathcal{R}_{\mathrm{w\text{-}dTNV}}(\bm{z})
    :=
    V_{\mathrm{vox}}\sum_{n=1}^{N} \phi\!\left(\tilde{\mathcal{G}}_n(\bm{z})\right),
    \qquad
    \phi(Y) := \sum_{\ell=1}^{\min(d,M)} g(\sigma_\ell(Y)),
    \label{eq:wdtnv_prior}
\end{equation}
where $V_{\mathrm{vox}}$ is the voxel volume and $g$ is the Charbonnier smoothing from~\eqref{eq:charbonnier_g}. Choosing equal weights $\rho_{m,n}\equiv 1$ recovers unweighted $d$TNV. Choosing $\rho_{m,n}$ proportional to the inverse of a characteristic modality scale brings both channels to a comparable range before the nuclear-norm coupling is evaluated. The \gls{pet}/\gls{spect} normalisation used in the application study is specified in Chapter~\ref{chapter:paper}.

\paragraph{Effect on the singular spectrum.}
Right-multiplying by $B_n$ scales the $m$-th column of $\tilde{\mathcal{G}}_n$ by $\rho_{m,n}$. This scales the singular values of $\tilde{\mathcal{G}}_n$ in a non-trivial way, since the singular values of $YB$ are not simply $\sigma_\ell(Y)\rho_m$. The intended effect is instead on the Gram matrix: $\tilde{\mathcal{G}}_n^\top \tilde{\mathcal{G}}_n$ contains modality self-terms and cross-terms after scaling, so suitable $\rho_{m,n}$ values allow both modalities to contribute comparably to the coupled singular spectrum. The $2\times2$ Gram-matrix algebra in Section~\ref{sec:tv_to_tnv} applies directly to this weighted case.

\section{Conclusion}

This chapter introduced a smoothed, weighted, and directional Total Nuclear Variation ($\wdtnv$) prior for synergistic reconstruction. Charbonnier smoothing of the singular values yields a differentiable spectral functional whose Jacobian-space gradient can be written in Gram-matrix form. For two coupled emission modalities this reduces the key spectral operation to a $2\times2$ inverse square-root, allowing the prior and its gradient to be evaluated robustly without a general per-voxel \glsentryshort{svd}.

The chapter also separated two roles that are easily confounded in multimodality reconstruction. Anatomical information is used directionally through $\Pi_{m,n}$, steering the penalty without adding \gls{ct} as another coupled intensity channel. Modality scaling is handled separately through the diagonal weight matrix $B_n$, allowing \gls{pet} and \gls{spect} gradients to enter the nuclear-norm coupling on comparable scales. Chapter~\ref{chapter:optimisation} uses these definitions to construct practical preconditioners and subset-selection strategies for large-scale $\wdtnv$ optimisation.
